\section{\centering ВВЕДЕНИЕ}
Численное решение задачи Коши для обыкновенных дифференциальных уравнений является одной из фундаментальных математических задач. Однако во многих случаях аналитическое решение таких задач невозможно \cite{Verzhbitsky}, поэтому возникла необходимость разработки и исследования различных математических моделей и численных методов. 

Наибольшее распространение среди методов получили методы Рунге-Кутты вследствие их высокой точности и универсальности. В реальных задачах часто применяются методы высоких порядков точности, построение которых связано с выполнением системы нелинейных условий на коэффициенты метода. Несмотря на то, что методы Рунге-Кутты изучаются давно, нахождение коэффициентов методов в явном виде затруднено вследствие существенного возрастания размера и сложности системы \cite{Demidovich}. 

Традиционные подходы к выбору коэффициентов, которые основаны на анализе локальных характеристик метода, являются недостаточными для ряда задач, в том числе для задач с осциллирующими решениями, вследствие того, что они не учитывают поведение численного решения на всём интервале интегрирования. Более того, в осциллирующих задачах стандартные методы демонстрируют численную ошибку, несообразную порядку точности \cite{article2023}. Поэтому возникает необходимость использования новых подходов для подбора коэффициентов.

Один из таких подходов заключается в рассмотрении семейства параметрических методов Рунге-Кутты, т.е. методов, содержащих свободные коэффициенты. Его преимущество заключается в возможности выбора из возможных схем наиболее подходящей для конкретной задачи, не фиксируя метод заранее. Именно оптимальный подбор параметров обеспечивает эффективность метода. Поэтому применение методов оптимизации для подбора численных схем является перспективным способом повышения точности численных методов за счёт адаптации метода к решаемой задаче \cite{article2023}. В частности, в ряде задач, где исследуются сложные многомерные пространства параметров, используется эволюционный алгоритм, преимущество которого заключается в отсутствии необходимости вычисления производных целевой функции \cite{article2025}.

Актуальность данной работы обусловлена необходимостью повышения точности численного решения задач Коши для обыкновенных дифференциальных уравнений с осциллирующими решениями, а также разработки методов, способных учитывать глобальные свойства решения на всём интервале интегрирования. Использование параметрических семейств методов Рунге-Кутты в сочетании с современными методами оптимизации позволяет существенно улучшить качество численного решения и расширить область применимости классических численных схем.

Целью данной работы является разработка и исследование параметрических методов Рунге–Кутты, ориентированных на эффективное решение задачи Коши с осциллирующими решениями, а также изучение подходов к оптимизации их коэффициентов.

Для достижения поставленной цели необходимо решить следующие задачи:
\begin{itemize}
	\item провести анализ существующих методов Рунге–Кутты и условий их построения;
	\item исследовать особенности численного решения задач с осциллирующими решениями;
	\item построить семейства параметрических методов Рунге–Кутты;
	\item разработать подходы к оптимизации коэффициентов методов, в том числе с использованием эволюционных алгоритмов;
	\item провести аналитическое исследование свойств полученных методов;
	\item реализовать разработанные методы и выполнить их численное тестирование;
	\item провести сравнительный анализ эффективности предложенных методов.
\end{itemize}

Таким образом, в работе рассматривается задача построения и исследования параметрических методов Рунге–Кутты, а также разработки эффективных алгоритмов оптимизации их коэффициентов для повышения точности численного решения задачи Коши.

 \clearpage
\section{\centering ОБЗОР ЛИТЕРАТУРЫ}

\subsection{Методы Рунге–Кутты}

\subsection{Задачи с осциллирующими решениями и параметрические методы}

\subsection{Нелинейное уравнение Шрёдингера и численные методы}

\subsection{Методы оптимизации коэффициентов}


 \clearpage
\section{\centering ПОСТАНОВКА ЗАДАЧИ}
%постановка задачи - уравнение, параментрические функции, надо вывести коэффициенты исходя из решения типовых задач

Рассматривается задача Коши для обыкновенного дифференциального уравнения, возникающая после пространственной дискретизации исходного уравнения в частных производных (например, нелинейного уравнения Шрёдингера):
\begin{equation}
	\frac{du}{dt} = f(t,u), \quad u(0) = u_0,
\end{equation}
где $u(t) \in \mathbb{C}^m$, $f: \mathbb{R} \times \mathbb{C}^m \to \mathbb{C}^m$.

Для численного решения задачи используется $s$-стадийный явный метод Рунге--Кутты:
\begin{equation}
	u_{n+1} = u_n + h \sum_{i=1}^{s} b_i k_i,
\end{equation}
\begin{equation}
	k_i = f\Bigl(t_n + c_i h,\; u_n + h \sum_{j=1}^{i-1} a_{ij} k_j \Bigr), \quad i = 1, \dots, s,
\end{equation}
где $h$ — шаг интегрирования, коэффициенты $a_{ij}, b_i, c_i$ определяют конкретный метод.

\medskip

Рассматривается параметрическое семейство методов Рунге--Кутты шестого порядка точности с $s = 8$ стадиями. Коэффициенты метода удовлетворяют системе нелинейных алгебраических уравнений (условиям порядка), обеспечивающих шестой порядок точности. Часть коэффициентов фиксируется, а оставшиеся рассматриваются как свободные параметры:
\begin{equation}
	\theta = (c_2, c_4, c_5, c_7) \in \mathbb{R}^4.
\end{equation}

Таким образом, формируется семейство методов
\begin{equation}
	\mathcal{M}(\theta),
\end{equation}
каждый из которых имеет одинаковый алгебраический порядок, но различные численные свойства.

\medskip

Целью работы является выбор таких значений параметров $\theta$, при которых метод $\mathcal{M}(\theta)$ обеспечивает наилучшее качество численного решения для задач с осциллирующими решениями.

Качество метода оценивается с помощью глобальной погрешности. Пусть $u(t_n)$ — точное решение, а $u_n$ — численное решение. Тогда глобальная погрешность определяется как
\begin{equation}
	E_n = \|u_n - u(t_n)\|.
\end{equation}

Для оценки метода используется максимум погрешности на интервале интегрирования:
\begin{equation}
	E = \max_{0 \le n \le N} \|u_n - u(t_n)\|.
\end{equation}

\medskip

Задача выбора оптимальных коэффициентов формулируется как задача многомерной нелинейной оптимизации. В частности, рассматривается целевая функция, характеризующая среднюю относительную глобальную ошибку по набору шагов интегрирования:
\begin{equation}
	J(\theta) = \frac{1}{N} \sum_{i=1}^{N} 
	\frac{\|E^{(i)}(\theta)\|_{\infty}}{\|\bar{E}^{(i)}\|_{\infty}},
\end{equation}
где:
\begin{itemize}
	\item $E^{(i)}(\theta)$ — глобальная ошибка исследуемого метода при шаге $h_i$,
	\item $\bar{E}^{(i)}$ — ошибка эталонного метода того же порядка,
	\item $N$ — число различных шагов интегрирования.
\end{itemize}

Таким образом, требуется решить задачу оптимизации:
\begin{equation}
	J(\theta) \to \min_{\theta \in \Omega},
\end{equation}
где $\Omega$ — допустимое множество параметров (например, $\theta \in [0,1]^4$).

\medskip

Дополнительно могут учитываться ограничения, связанные со свойствами метода, такими как устойчивость и величина коэффициентов, например:
\begin{enumerate}
	\item обеспечение достаточного интервала устойчивости;
	\item ограниченность коэффициентов для снижения накопления округлённых ошибок.
\end{enumerate}

\medskip


\clearpage
\section{\centering РАЗРАБОТКА МЕТОДОВ}


%как применить метод оптимизации

\clearpage
\section{\centering РЕЗУЛЬТАТЫ}

\clearpage
\subsection{Сравнение с решением, полученным машинным обучением}

\clearpage
\section{ВЫВОДЫ}

\clearpage
\section{ЗАКЛЮЧЕНИЕ}
%цели, выводы, заключение - краткий пересказ, зеркально к введению

% введение - актуальность, research gap, обзор поделить на методы рунге-кутты (страница - полторы), потом новые задачи и методы, обзор по методам оптимизации

%описать уравнения шредингера в обзоре



% -> плавно переходим к методам, которые разрабатываются


% начинать с самого метода интегрирования - прогоняь интегрированное уравнение с разными коэффициентами

% как начать?

% получается несколько модулей которые и сами по себе что-то делают и должны между собой контактировать

% 1) метод рунге-кутты умеет интегрировать разные уравнения (шредингера, эн боди, и тд). Это метод 8 порядка, какие-то коэффициенты фиксированы, а какие-то можно менять. Это файл где нам интересна одна функция, анализ метода. Принимает правую часть уравнения, то что надо в задаче Коши (н.у., временной промежуток, начальные значения шага и тд) и плавающие параметры. Через класс, списки, как угодно. Выдавать должна таблицу, после того как проинтегрировала

% 2) второй модуль - должен непосредственно уравнения шредингера и тд

% 3) еще модуль - считает целевую функцию - перебирает h, вычисляет одним методом, другим, вычитает и получает какую-то ошибку. Этот модуль должен вызывать, должны быть плавающие параметры которые ищем

% модули, ограничения, может посчитать значения

% 5) когда это будет работать можно будет оптимизировать - по большому счету два варианта: та которую они использовали и надо получить что-то похожее. Разностная эволюция